\begin{table}[!t]
\centering
\caption{Descriptive Statistics: Food and Beverage Service Enterprises}
\label{tab:descriptives}
\begin{tabular}{lcccc}
\hline\hline
 & Total & Large (250+) & Large & Threshold \\
 & Enterprises & Enterprises & Share (\%) & Ratio \\
\hline
\multicolumn{5}{l}{\textit{Panel A: England (Treated)}} \\[3pt]
Pre-regulation (2010--2022) & 109,071 & 348 & 0.317 & 0.409 \\
Post-regulation (2023--2024) & 129,705 & 442 & 0.341 & 0.428 \\[6pt]
\multicolumn{5}{l}{\textit{Panel B: Scotland (Control)}} \\[3pt]
Pre-regulation (2010--2022) & 10,608 & 13 & 0.120 & 0.204 \\
Post-regulation (2023--2024) & 12,028 & 20 & 0.166 & 0.310 \\[6pt]
\multicolumn{5}{l}{\textit{Panel C: Wales (Control)}} \\[3pt]
Pre-regulation (2010--2022) & 6,808 & 8 & 0.110 & 0.500 \\
Post-regulation (2023--2024) & 7,645 & 15 & 0.196 & 0.333 \\
\hline\hline
\end{tabular}
\begin{tablenotes}[flushleft]
\small
\item \textit{Notes:} Data from ONS UK Business Counts (NOMIS NM\_142\_1), 
annual March snapshots, 2010--2024. ``Total Enterprises'' counts all registered 
food and beverage service enterprises (SIC 56). ``Large (250+)'' counts enterprises 
with 250 or more employees, subject to England's mandatory calorie labeling regulation 
(effective April 6, 2022). ``Large Share'' is the ratio of large to total enterprises. 
``Threshold Ratio'' is enterprises with 250--499 employees divided by those with 
100--249 employees. Scotland and Wales did not adopt calorie labeling mandates. 
Enterprise counts are rounded to the nearest 5 by ONS.
\end{tablenotes}
\end{table}
